% Options for packages loaded elsewhere
\PassOptionsToPackage{unicode}{hyperref}
\PassOptionsToPackage{hyphens}{url}
\documentclass[
]{article}
\usepackage{xcolor}
\usepackage{amsmath,amssymb}
\setcounter{secnumdepth}{-\maxdimen} % remove section numbering
\usepackage{iftex}
\ifPDFTeX
  \usepackage[T1]{fontenc}
  \usepackage[utf8]{inputenc}
  \usepackage{textcomp} % provide euro and other symbols
\else % if luatex or xetex
  \usepackage{unicode-math} % this also loads fontspec
  \defaultfontfeatures{Scale=MatchLowercase}
  \defaultfontfeatures[\rmfamily]{Ligatures=TeX,Scale=1}
\fi
\usepackage{lmodern}
\ifPDFTeX\else
  % xetex/luatex font selection
\fi
% Use upquote if available, for straight quotes in verbatim environments
\IfFileExists{upquote.sty}{\usepackage{upquote}}{}
\IfFileExists{microtype.sty}{% use microtype if available
  \usepackage[]{microtype}
  \UseMicrotypeSet[protrusion]{basicmath} % disable protrusion for tt fonts
}{}
\makeatletter
\@ifundefined{KOMAClassName}{% if non-KOMA class
  \IfFileExists{parskip.sty}{%
    \usepackage{parskip}
  }{% else
    \setlength{\parindent}{0pt}
    \setlength{\parskip}{6pt plus 2pt minus 1pt}}
}{% if KOMA class
  \KOMAoptions{parskip=half}}
\makeatother
\setlength{\emergencystretch}{3em} % prevent overfull lines
\providecommand{\tightlist}{%
  \setlength{\itemsep}{0pt}\setlength{\parskip}{0pt}}
\usepackage{bookmark}
\IfFileExists{xurl.sty}{\usepackage{xurl}}{} % add URL line breaks if available
\urlstyle{same}
\hypersetup{
  hidelinks,
  pdfcreator={LaTeX via pandoc}}

\author{}
\date{}

\begin{document}

\section{Independent check of
BW13--24}\label{independent-check-of-bw1324}

Read ROOT\_BOUNDARY\_WEIGHT\_HOMOLOGY.tex in full; the bounded
mathematical review concerns BW13--24. The finite Koszul contraction,
balanced transition maps, all-prime homology colimit, original idele
action, stated multiplicities, and labelled fibers are correct.

The proof of the Schwartz-boundary connecting map needs an explicit
sign-subgroup map. The element −1 acts trivially on the boundary module
V\_r, but it generally acts nontrivially on the original Schwartz test
module T\_r and its kernel J\_r.

Let A=U\_\{−1\}\^{}\{tensor r\}, so A²=I, and let P\_+=(I+A)/2. For
either original module M, write M\^{}+=P\_+M. The decomposition M=M\^{}+
direct-sum M\^{}- is exact, and passing to the sign subgroup's
coinvariants is canonically the map P\_+. Its kernel is M\^{}-: for v in
M\^{}-, v=(A−I)(−v/2). The averaging projector commutes with every
rational and idele dilation.

The full group is the direct product of the order-two sign group and the
free abelian positive-rational group. Over C, the trivial module of the
sign group is projective, by its averaging idempotent. Its degree-zero
projective resolution, tensored with the original positive-prime Koszul
resolution, therefore computes the full group homology using the
coefficient modules M\^{}+.

Applying P\_+ to BW19 gives the exact original sequence

\begin{verbatim}
0 -> J_r^+ -> T_r^+ -> V_r -> 0.
\end{verbatim}

In particular, an actual lift t of any boundary coefficient can and
should be replaced in the full-group connecting calculation by
t\^{}+=(t+At)/2. Its boundary coefficient is unchanged because the
boundary action of −1 is trivial. The Koszul differential of this
averaged lift lies in J\_r\^{}+ and gives the stated full-group
connecting class. Changing the lift changes that class by a kernel
differential, exactly as required.

BW22 remains valid with any unaveraged test lift. If u=(U\_p\^{}\{tensor
r\}−I)t belongs to J\_r, its averaged version is P\_+u. The difference
u−P\_+u is in the negative sign eigenspace and equals
(A−I)(−(u−P\_+u)/2), so it has zero class in (J\_r)\_\{Q×\}. Thus the
original coefficient formula is preserved; the missing averaging map
completes its full-group justification.

For the infinite-resolution paragraph, an additional explicit algebraic
sentence makes the augmentation proof complete: the all-prime Laurent
ring R is free over each R\_F on the Laurent monomials in the other
variables. Hence base change preserves the finite Koszul exact sequence
and resolves R/(t\_p−1:p in F). The filtered colimit of these quotient
modules is C. The union of the free Koszul modules is R tensor Lambda E,
with the actual ordered prime basis, giving the asserted free
resolution.

The labelled maps BW23--24 preserve the entire support lattice. They are
lifts of linear maps, not a claim that amplitude exactness deletes a
support label. A full-group sign averaging map itself also has this
same-label lift; its amplitude factor 1/2 does not change the support
coordinate.

The only additional source typo reported was the missing backslash
before quad in BW11. No new Frobenius, interior-cohomology, or
Riemann-hypothesis conclusion follows from this bounded check.

\end{document}
